\chapter{Offene Fragen und ein Forschungsprogramm}
\label{ch:offen}

% Quelle: conjectures.md, Stand Version 0.7.

Bevor das letzte Kapitel die Theorie auf der Landkarte der
Mathematik verortet, ein ehrlicher Zwischenhalt: nicht, was
bewiesen wurde -- das haben die vorigen Kapitel getan --, sondern
was offen ist, und in welcher Reihenfolge es sich zu versuchen
lohnt.

\section{Die nächstliegenden Fragen}

\paragraph{Die Ziffernbilder der Fakultäten.} Die Ziffernformel für
die Schalenspringer hat die Frage verschoben: Nicht mehr „wie viele
Springer gibt es?“, sondern „wie schreibt sich $n!$ in seiner
eigenen Schale?“. Erste Muster liegen vor: Die Zahl der führenden
Nullen wächst wie die Umkehrfunktion der Fakultät, und die
Ein-Ziffern-Fälle sind inzwischen \emph{vollständig}
charakterisiert: Es sind exakt die $n=(i+1)!/d-1$ mit $1\le d\le i$,
genau $i$ Stück pro Position, und die Vorhersagen $n=29,39,59,119$
sowie $143,179,239,359,719$ trafen alle ein
(Anhang~\ref{ch:anhang-beweise}). Eine vollständige Theorie der
übrigen Ziffernbilder (Verteilung der Nichtnull-Anzahl, Position der
ersten Maximalziffer) fehlt weiterhin.

\paragraph{Die Feinstruktur der Interchange-Gleichheit.} Die
Positivität des Interchange-Defekts und die Charakterisierung der
Gleichheitsfälle sind bewiesen (Anhang~\ref{ch:anhang-beweise});
die Zählfolge $2$, $5$, $13$, $23$, $61$, $111$, $274$,
$564,\ldots$ gehorcht einer semi-geschlossenen Formel über die
scharfen Quotienten $q'$ aus dem Positivitätsbeweis, und
deren vollständige Ziffernentfaltung ist inzwischen bewiesen: pro
Ebene zwei Ungleichungen, ein wachsender Exzess, ein Abstieg
(Anhang~\ref{ch:anhang-beweise}). Die Asymptotik des
Nullzähler-Anteils ist inzwischen nach oben geklärt: Gleichgültigkeit gegenüber der Wachstumsreihenfolge ist
asymptotisch extrem selten -- der Anteil fällt mindestens wie
$e^{-(1+o(1))\,n\ln n}$, weil die Scharfheit die Ziffern ebenenweise
abwürgt. Offen bleibt die passende untere Schranke.

\paragraph{Die obere Seite des Tiefengesetzes.} Die untere Seite
ist inzwischen vollständig bewiesen -- Sprungstellen, Asymptotik
und, als letztes Stück während der Schlussredaktion, die exakte
Gleichheit: Die maximale Tiefe in Schale $m$ ist \emph{genau} die
Schalenlücke, und ein stets tauglicher Zeuge ist die
Terminaldarstellung des Maximalpolynom-Wertes (für Schale 3: die
Hausgeist-Zahl $(6,6)$).
Offen ist nur noch die obere Seite: Wie schnell wächst das
\emph{Maximum} des Kommutators? Hier ist Vorsicht geboten, und zwar
methodisch begründete: Exhaustiv gesichert ist nur der erste Sprung
($+2$ ab der neunten Welt). Alle späteren Schwellen sind bloß nach
oben beschränkt ($+3$ spätestens ab der vierzehnten, $+4$ spätestens
ab der zwanzigsten, $+5$ spätestens ab der sechsundzwanzigsten) --
jede Vertiefung der Zeugensuche hat sie bisher nach vorn verschoben,
und zwei voreilige Formelhypothesen sind an genau dieser Front
gestorben, eine davon noch vor ihrer Niederschrift.

\section{Die strukturellen Fragen}

\paragraph{Gekoppelte Arithmetiken.} Der Seitentreue-Satz verbietet
das Mischen der Identitäten; \emph{innerhalb} einer Seite ist die
Landschaft aber nicht klassifiziert. Die Produktkonstruktionen über
den Exzess liefern eine ganze Familie; der konservativ-kompakte
Halbring ist nachweislich keine davon. Gibt es weitere gekoppelte
Strukturen? Gibt es eine vollständige Klassifikation pro Seite?

\paragraph{Nachtrag: der nichtkommutative Mischfall ist geschlossen.}
Auch diese Lücke fiel während der Schlussredaktion: Verlangt man,
wie in jedem Halbring, \emph{beide} Distributivgesetze, so spiegelt
die Rechts-Distributivität die Kettenkonstruktion und ersetzt die
Kommutativität vollständig (Beweis in
Anhang~\ref{ch:anhang-beweise}). Was übrig bleibt, ist nur die
Kuriosität eines einseitig distributiven Mischsystems -- außerhalb
jeder Halbring-Axiomatik, und ehrlicherweise eher eine Fußnote als
eine Frage.

\paragraph{Nachtrag: die dreiseitige Seitentreue ist bewiesen.}
Während der Schlussredaktion fiel auch der letzte Fall
(Gestalt-Addition mit Bruch-Multiplikation): Die Kettengleichung
erzwingt die polynomiale Identität $2A(Y)=Y^{\underline{k}}$, die
den Koeffizienten $\tfrac12$ verlangt -- Ziffern sind aber ganz.
Zusammen mit den zwei zuvor entschiedenen Fällen gilt: Es gibt genau
zwei Rechenwelten (Wert und Gestalt), und der Bruch trägt keine
dritte und mischt sich in keine. Details in
Anhang~\ref{ch:anhang-beweise}. Offen bleibt hier also -- nichts
mehr; die Frage wandert als bewiesener Satz in
Kapitel~\ref{ch:arithmetik}.

\paragraph{Formalisierung.} Der Seitentreue-Satz ist elementar
genug für eine vollständige Formalisierung in Lean; die fallenden
Fakultäten liegen in Mathlib bereits vor. Das wäre der Goldstandard
der Absicherung -- und ein Beitrag zur Formal-Methods-Gemeinde
gleich mit.

\section{Woran sich die Theorie messen lassen will}

Zum Schluss die Frage hinter allen Fragen: Wann wäre die Horimetrik
\emph{gescheitert}, wann \emph{gelungen}? Gescheitert, wenn sich
ein bestehender Formalismus findet, in dem Horizont, Dualität und
Seitentreue nur Umbenennungen sind -- die Literatursuche hat ihn
nicht gefunden, aber Suchen beweisen nichts. Gelungen, wenn eines
von dreien eintritt: wenn die Zählfolgen dieses Buchs anderswo
unabhängig auftauchen; wenn der Seitentreue-Satz oder das
Tiefengesetz in fremder Sprache neu bewiesen wird; oder wenn
jemand, der dieses Buch nie gelesen hat, an den Ziffern von $n!$
in der eigenen Schale dieselbe Freude findet wie wir.

\begin{oberflaeche}
Eine Theorie ist fertig, wenn sie keine Fragen mehr erzeugt. Nach
diesem Maßstab ist die Horimetrik erfreulich unfertig -- und genau
so soll dieser Zwischenhalt enden: nicht mit einem Punkt, sondern
mit einem Doppelpunkt. Was bleibt, ist die Frage nach der Adresse:
Wo auf der Landkarte wohnt das alles?
\end{oberflaeche}
